\documentclass[aspectratio=169,professionalfonts]{beamer}

\usepackage[T1]{fontenc}
\usepackage[utf8]{inputenc}
\usepackage[ngerman]{babel}
\usepackage{lmodern}
\usepackage{hyperref}
\usepackage{tikz}
\usepackage{xcolor}
\usetikzlibrary{decorations.pathreplacing}

\definecolor{wsdblue}{RGB}{18,53,91}
\definecolor{wsdgray}{RGB}{88,94,102}
\definecolor{wsdlight}{RGB}{240,243,246}

\setbeamertemplate{navigation symbols}{}
\setbeamertemplate{footline}{%
  \leavevmode%
  \hbox{%
    \begin{beamercolorbox}[wd=\paperwidth,ht=2.8ex,dp=1.2ex,leftskip=1em,rightskip=1em]{author in head/foot}%
      \color{wsdgray}\scriptsize Christopher Albert \hfill Wahrscheinlichkeitstheorie, Statistik und Datenanalyse \hfill \insertframenumber
    \end{beamercolorbox}%
  }%
}
\setbeamertemplate{headline}{}
\setbeamercolor{background canvas}{bg=white}
\setbeamercolor{normal text}{fg=black,bg=white}
\setbeamercolor{structure}{fg=wsdblue}
\setbeamercolor{frametitle}{fg=wsdblue,bg=white}
\setbeamerfont{title}{series=\bfseries,size=\LARGE}
\setbeamerfont{frametitle}{series=\bfseries,size=\Large}
\setbeamerfont{subtitle}{size=\large}

\title{Wahrscheinlichkeitstheorie, Statistik und Datenanalyse}
\author{Christopher Albert}
\date{26. M\"arz 2026}

\begin{document}

\begin{frame}[plain]
  \vspace{1.2cm}
  {\usebeamerfont{title}\color{wsdblue}\inserttitle\par}
  \vspace{0.5cm}
  {\usebeamerfont{subtitle}VO 2: Zufallsvariable und Mittelwert\par}
  \vspace{0.3cm}
  {\large Christopher Albert\par}
  \vfill
  {\large 26. M\"arz 2026\par}
\end{frame}

\begin{frame}
  \frametitle{R\"uckblick auf VO 1: Grundbegriffe}

  \begin{itemize}
    \item \textbf{Versuch}: ein einzelner Zufallsvorgang, z.B. einmal w\"urfeln
    \item \textbf{Ergebnis} \(\omega\): ein einzelner m\"oglicher Ausgang, z.B. die Zahl \(3\)
    \item \textbf{Ergebnisraum} \(\mathcal G\): die Menge aller m\"oglichen Ergebnisse, z.B.
    \[
      \mathcal G=\{1,2,3,4,5,6\}
    \]
    \item \textbf{Ereignis}: eine Teilmenge von \(\mathcal G\), z.B. \(\{2,4,6\}\)
    \item \textbf{Elementarereignis}: ein Ereignis mit genau einem Ergebnis, also \(\{\omega\}\), z.B. \(\{3\}\)
  \end{itemize}
\end{frame}

\begin{frame}
  \frametitle{Von Ergebnissen zu Zahlen}

  \begin{columns}[T,onlytextwidth]
    \column{0.48\textwidth}
    \textbf{W\"urfel}
    \begin{itemize}
      \item Ergebnisse: \(1,2,3,4,5,6\)
      \item hier sind die Ergebnisse schon Zahlen
      \item Ereignisse sind Teilmengen, z.B.
      \[
        \{1,2\},\qquad \{2,4,6\}
      \]
    \end{itemize}

    \column{0.48\textwidth}
    \textbf{Karte ziehen}
    \begin{itemize}
      \item Ergebnisse: konkrete Karten
      \item zun\"achst keine nat\"urlichen Zahlen
      \item zum Rechnen braucht man eine Abbildung auf reelle Zahlen
    \end{itemize}
  \end{columns}
\end{frame}

\begin{frame}
  \frametitle{Zufallsvariable}

  \begin{block}{Definition}
    Eine Zufallsvariable ist eine Abbildung
    \[
      X:\mathcal G\to \mathbb R,\qquad \omega \mapsto x=X(\omega)
    \]
    die jedem \textbf{Ergebnis} \(\omega\in\mathcal G\) eine reelle Zahl zuordnet.
  \end{block}

  \vspace{0.3cm}

  \begin{itemize}
    \item \(x\): Realisierung von \(X\)
    \item \(R\): Wertebereich der möglichen Realisierungen
    \item wichtig: \textbf{nicht} jedem Ereignis, sondern jedem Ergebnis
  \end{itemize}
\end{frame}

\begin{frame}
  \frametitle{Beispiel 1: Identische Abbildung}

  \[
    X(\omega)=\omega
  \]

  \begin{center}
    \input{skriptum/Figures/fig_random_variable_identity.tex}
  \end{center}

  \vspace{0.1cm}

  \begin{itemize}
    \item beim Würfel ist \(X(\omega)=\omega\) die naheliegende Wahl
  \end{itemize}
\end{frame}

\begin{frame}
  \frametitle{Beispiel 2: Indikatorfunktion}

  \[
    I_{A_{\mathrm{gerade}}}(\omega)=
    \begin{cases}
      1,& \omega\in\{2,4,6\}\\
      0,& \text{sonst}
    \end{cases}
  \]

  \begin{center}
    \input{skriptum/Figures/fig_random_variable_indicator.tex}
  \end{center}

  \vspace{0.2cm}

  \begin{itemize}
    \item \(I_{A_{\mathrm{gerade}}}(\omega)\) nimmt nur die Werte \(0\) und \(1\) an
    \item so markiert die Zufallsvariable genau das Ereignis \(A_{\mathrm{gerade}}=\{2,4,6\}\)
  \end{itemize}
\end{frame}

\begin{frame}
  \frametitle{Ereignisse \"uber eine Zufallsvariable definieren}

  \[
    A_{\mathrm{gerade}}=\{\omega\in\mathcal G:\,I_{A_{\mathrm{gerade}}}(\omega)>0.5\}
  \]

  \begin{center}
    \input{skriptum/Figures/fig_random_variable_levelset.tex}
  \end{center}

  \vspace{0.2cm}

  \begin{itemize}
    \item allgemeiner:
    \[
      \{\omega\in\mathcal G:\,X(\omega)\in B\}
    \]
    \item Ereignisse erscheinen als Urbilder von Wertebereichen
    \item das ist abstrakter und sp\"ater viel flexibler
  \end{itemize}
\end{frame}

\begin{frame}
  \frametitle{R\"uckblick: Stichprobe im Produktraum}

  \[
    \mathcal G^N=\underbrace{\mathcal G\times\cdots\times\mathcal G}_{N\text{-mal}}
  \]

  \[
    \boldsymbol{\omega}=(\omega_1,\ldots,\omega_N)\in\mathcal G^N
  \]

  \begin{itemize}
    \item das haben wir bei der Wiederholung eines Versuchs schon in VO1 eingef\"uhrt
    \item eine Stichprobe vom Umfang \(N\) ist ein geordnetes \(N\)-Tupel von Ergebnissen
    \item darauf definieren wir jetzt Statistiken wie den Stichprobenmittelwert
  \end{itemize}
\end{frame}

\begin{frame}
  \frametitle{Beispiel: Dreimal Münzwurf}

  \[
    \mathcal G^3=\{K,Z\}^3
  \]

  \begin{center}
    \begin{tabular}{c c}
      Ergebnis & Zahl der K\"opfe \\
      \hline
      \((K,K,K)\) & \(3\) \\
      \((K,K,Z)\) & \(2\) \\
      \((K,Z,K)\) & \(2\) \\
      \((Z,K,K)\) & \(2\) \\
      \((K,Z,Z)\) & \(1\) \\
      \((Z,K,Z)\) & \(1\) \\
      \((Z,Z,K)\) & \(1\) \\
      \((Z,Z,Z)\) & \(0\)
    \end{tabular}
  \end{center}

  \vspace{0.2cm}

  \[
    P(\omega_1,\omega_2,\omega_3)=P_{\omega_1}\,P_{\omega_2}\,P_{\omega_3}
  \]
\end{frame}

\begin{frame}
  \frametitle{Funktionen von Zufallsvariablen}

  \begin{itemize}
    \item aus einer Zufallsvariablen \(X\) kann man wieder neue machen:
    \[
      Y=f(X)
    \]
    \item jedem Ergebnis \(\omega\) wird dann der Wert \(f(X(\omega))\) zugeordnet
    \item damit lassen sich Gewinne, Schwellen, Indikatoren oder Klassifikationen beschreiben
  \end{itemize}
\end{frame}

\begin{frame}
  \frametitle{Mittelwert einer diskreten Zufallsvariablen}

  \[
    \langle X\rangle=\sum_{\omega\in\mathcal G} X(\omega)\,P_\omega
  \]

  \vspace{0.35cm}

  \begin{itemize}
    \item gewichtete Summe der Werte mit ihren Wahrscheinlichkeiten
    \item oft auch \textbf{Erwartungswert} genannt
    \item wir sagen hier bevorzugt \textbf{Mittelwert}
  \end{itemize}
\end{frame}

\begin{frame}[plain]
\end{frame}

\begin{frame}
  \frametitle{Fairer W\"urfel: Was ist der Mittelwert?}

  \[
    X(\omega)=\omega,\qquad P_n=\frac{1}{6}
  \]

  \vspace{0.35cm}

  \[
    \langle X\rangle
    =\frac{1}{6}\sum_{n=1}^{6} n
    =\frac{1}{6}\cdot\frac{6\cdot 7}{2}
    =3.5
  \]

  \vspace{0.35cm}

  \begin{itemize}
    \item genau deshalb ist \glqq Mittelwert\grqq\ hier das bessere Wort
    \item niemand erwartet, dass der W\"urfel jemals \(3.5\) zeigt
  \end{itemize}
\end{frame}

\begin{frame}
  \frametitle{Mittelwert versus Stichprobenmittelwert}

  \begin{columns}[T,onlytextwidth]
    \column{0.48\textwidth}
    \textbf{Mittelwert \(\langle X\rangle\)}
    \begin{itemize}
      \item fixer Kennwert der Verteilung
      \item keine Zufallsvariable
      \item auch: Erwartungswert
    \end{itemize}

    \column{0.48\textwidth}
    \textbf{Stichprobenmittelwert \(\overline X\)}
    \begin{itemize}
      \item aus beobachteten Daten berechnet
      \item Statistik auf dem Produktraum \(\mathcal G^N\)
      \item Sch\"atzer für \(\langle X\rangle\)
    \end{itemize}
  \end{columns}

  \vspace{0.3cm}

  \[
    \overline X(\omega_1,\ldots,\omega_N)=\frac{1}{N}\sum_{i=1}^{N}X(\omega_i)
  \]
\end{frame}

\begin{frame}
  \frametitle{Linearit\"at des Mittelwerts}

  \[
    \langle \alpha f(X)+\beta g(X)\rangle
    =
    \alpha \langle f(X)\rangle+\beta \langle g(X)\rangle
  \]

  \vspace{0.4cm}

  \begin{itemize}
    \item Konstanten kann man herausziehen
    \item Summen kann man im Mittel trennen
    \item wenn alles doppelt so gro\ss\ ist, ist es im Mittel auch doppelt so gro\ss
  \end{itemize}
\end{frame}

\begin{frame}
  \frametitle{Momente}

  \[
    m_i=\langle X^i\rangle
  \]

  \vspace{0.3cm}

  \begin{itemize}
    \item das \(i\)-te Moment ist der Mittelwert der Potenz \(X^i\)
    \item \(m_0=1\) wegen der Normierung
    \item \(m_1=\langle X\rangle\), \quad \(m_2=\langle X^2\rangle\), \quad \(m_3=\langle X^3\rangle\)
    \item Momente charakterisieren die Verteilung feiner als nur ihr Mittelwert
  \end{itemize}
\end{frame}

\begin{frame}
  \frametitle{Zentrale Momente und Varianz}

  \[
    \mu_i=\langle (X-\langle X\rangle)^i\rangle
  \]

  \[
    \mathrm{Var}(X)=\sigma_X^2
    =
    \langle (X-\langle X\rangle)^2\rangle
    =
    \langle X^2\rangle-\langle X\rangle^2
  \]

  \vspace{0.25cm}

  \begin{itemize}
    \item zentrale Momente messen Abweichungen um den Mittelwert
    \item die Varianz ist das zweite zentrale Moment
    \item \glqq Verschiebungssatz\grqq: zweites nichtzentrales Moment minus Quadrat des Mittelwerts
  \end{itemize}
\end{frame}

\begin{frame}[plain]
\end{frame}

\begin{frame}
  \frametitle{Fairer W\"urfel: Varianz und Standardabweichung}

  \[
    \langle X^2\rangle=\frac{1}{6}\sum_{n=1}^{6} n^2=\frac{91}{6}\approx 15.17
  \]

  \[
    \mathrm{Var}(X)=\langle X^2\rangle-\langle X\rangle^2
    =\frac{91}{6}-3.5^2
    =\frac{35}{12}\approx 2.92
  \]

  \[
    \sigma_X=\sqrt{\mathrm{Var}(X)}=\sqrt{\frac{35}{12}}\approx 1.71
  \]
\end{frame}

\begin{frame}
  \frametitle{Kennzahl der Verteilung oder Statistik?}

  \begin{columns}[T,onlytextwidth]
    \column{0.48\textwidth}
    \textbf{Kennzahlen der Verteilung}
    \begin{itemize}
      \item Mittelwert
      \item Momente
      \item Varianz
      \item Standardabweichung
    \end{itemize}

    \column{0.48\textwidth}
    \textbf{Statistiken aus Daten}
    \begin{itemize}
      \item Stichprobenmittelwert
      \item Stichprobenvarianz
      \item h\"angen von der konkreten Stichprobe ab
      \item n\"ahern sich bei gro\ss em \(N\) den wahren Kennzahlen an
    \end{itemize}
  \end{columns}
\end{frame}

\begin{frame}
  \frametitle{Standardfehler des Stichprobenmittelwerts}

  \[
    \mathrm{SE}(\overline X)=\frac{\sigma_X}{\sqrt N}
  \]

  \vspace{0.3cm}

  \begin{itemize}
    \item misst die Streuung des Stichprobenmittelwerts, nicht der Einzelwerte
    \item f\"allt nur wie \(1/\sqrt N\)
    \item mehr Daten helfen, aber nicht linear
    \item sp\"ater verbinden wir das mit dem Gesetz der gro\ss en Zahlen
  \end{itemize}
\end{frame}

\begin{frame}
  \frametitle{Stichprobenvarianz}

  \[
    s^2=\frac{1}{N-1}\sum_{i=1}^{N}(x_i-\overline x)^2
  \]

  \vspace{0.3cm}

  \begin{itemize}
    \item wichtig: \(N-1\), nicht \(N\)
    \item sonst untersch\"atzt man die Streuung systematisch
    \item Grund: der Mittelwert \(\overline x\) wurde aus denselben Daten schon gesch\"atzt
    \item bei kleinen Stichproben ist dieser Unterschied besonders relevant
  \end{itemize}
\end{frame}

\begin{frame}
  \frametitle{Bild: Gleicher Mittelwert, andere Varianz}

  \[
    \langle X\rangle = 3.5 \qquad \text{in beiden Verteilungen}
  \]

  \begin{center}
    \begin{tikzpicture}[x=1.15cm,y=0.58cm]
      \draw[red!75!black,dashed,thick] (3.5,0.5) -- (3.5,7.0);
      \node[red!75!black,anchor=south,font=\small] at (3.5,7.0) {Mittelwert \(3.5\)};

      \node[anchor=east,font=\small\bfseries] at (0.35,5.55) {enger};
      \draw[thick] (0.7,4.7) -- (6.3,4.7);
      \foreach \x in {1,...,6} {
        \draw[thick] (\x,4.58) -- (\x,4.82);
        \node[font=\scriptsize] at (\x,4.2) {\x};
      }
      \foreach \y in {5.1,5.45,5.8,6.15} {
        \fill[wsdblue] (3,\y) circle (2.6pt);
        \fill[wsdblue] (4,\y) circle (2.6pt);
      }
      \node[anchor=west,font=\small] at (6.45,5.55) {\(P(3)=P(4)=\frac12\)};
      \node[anchor=west,font=\small] at (6.45,5.0) {\(\mathrm{Var}(X)=0.25\)};

      \node[anchor=east,font=\small\bfseries] at (0.35,2.25) {breiter};
      \draw[thick] (0.7,1.4) -- (6.3,1.4);
      \foreach \x in {1,...,6} {
        \draw[thick] (\x,1.28) -- (\x,1.52);
        \node[font=\scriptsize] at (\x,0.9) {\x};
      }
      \foreach \y in {1.8,2.15,2.5,2.85} {
        \fill[wsdblue] (1,\y) circle (2.6pt);
        \fill[wsdblue] (6,\y) circle (2.6pt);
      }
      \node[anchor=west,font=\small] at (6.45,2.25) {\(P(1)=P(6)=\frac12\)};
      \node[anchor=west,font=\small] at (6.45,1.7) {\(\mathrm{Var}(X)=6.25\)};
    \end{tikzpicture}
  \end{center}

  \vspace{0.15cm}

  \begin{itemize}
    \item der Mittelwert markiert die Lage des Schwerpunkts
    \item die Varianz misst, wie weit die Wahrscheinlichkeitsmasse um diesen Schwerpunkt verteilt ist
  \end{itemize}
\end{frame}

\begin{frame}
  \frametitle{Bild: Gleicher Stichprobenmittelwert, andere Stichprobenvarianz}

  \[
    \overline x = 3.5 \qquad \text{in beiden Stichproben}
  \]

  \begin{center}
    \begin{tikzpicture}[x=1.15cm,y=0.58cm]
      \draw[red!75!black,dashed,thick] (3.5,0.5) -- (3.5,6.4);
      \node[red!75!black,anchor=south,font=\small] at (3.5,6.4) {Stichprobenmittelwert \(3.5\)};

      \node[anchor=east,font=\small\bfseries] at (0.35,5.15) {Stichprobe A};
      \draw[thick] (0.7,4.3) -- (6.3,4.3);
      \foreach \x in {1,...,6} {
        \draw[thick] (\x,4.18) -- (\x,4.42);
        \node[font=\scriptsize] at (\x,3.8) {\x};
      }
      \fill[wsdblue] (3,4.7) circle (2.8pt);
      \fill[wsdblue] (3,5.05) circle (2.8pt);
      \fill[wsdblue] (4,4.7) circle (2.8pt);
      \fill[wsdblue] (4,5.05) circle (2.8pt);
      \node[anchor=west,font=\small] at (6.45,5.15) {\(3,3,4,4\)};
      \node[anchor=west,font=\small] at (6.45,4.6) {\(s^2=\frac13\)};

      \node[anchor=east,font=\small\bfseries] at (0.35,2.15) {Stichprobe B};
      \draw[thick] (0.7,1.3) -- (6.3,1.3);
      \foreach \x in {1,...,6} {
        \draw[thick] (\x,1.18) -- (\x,1.42);
        \node[font=\scriptsize] at (\x,0.8) {\x};
      }
      \fill[wsdblue] (1,1.7) circle (2.8pt);
      \fill[wsdblue] (2,1.7) circle (2.8pt);
      \fill[wsdblue] (5,1.7) circle (2.8pt);
      \fill[wsdblue] (6,1.7) circle (2.8pt);
      \node[anchor=west,font=\small] at (6.45,2.15) {\(1,2,5,6\)};
      \node[anchor=west,font=\small] at (6.45,1.6) {\(s^2=\frac{17}{3}\)};
    \end{tikzpicture}
  \end{center}

  \vspace{0.15cm}

  \begin{itemize}
    \item der Stichprobenmittelwert beschreibt die Lage der beobachteten Daten
    \item die Stichprobenvarianz beschreibt ihre Streuung um diesen gesch\"atzten Mittelpunkt
  \end{itemize}
\end{frame}

\begin{frame}
  \frametitle{Bis Hierher}

  \begin{itemize}
    \item Zufallsvariablen bilden Ergebnisse auf reelle Zahlen ab
    \item Ereignisse kann man \"uber Zufallsvariablen als Urbilder definieren
    \item bei Wiederholung entsteht nat\"urlich ein Produktraum
    \item Mittelwert, Momente, Varianz und Standardabweichung sind Kennzahlen der Verteilung
    \item Stichprobenmittelwert und Stichprobenvarianz sind Statistiken aus endlichen Daten
    \item gleicher Mittelwert hei\ss t noch nicht gleiche Streuung
    \item der Standardfehler beschreibt die Unsicherheit des Stichprobenmittelwerts
  \end{itemize}
\end{frame}

\end{document}
